
\lussec 2. The gradient flow in QCD

The theory considered in this paper is QCD with gauge group $\SUn$ and a 
multiplet of two or more massive quarks in the 
fundamental representation of the gauge group. 
Many results are however of a fairly general nature
and not limited to QCD. The theory is set up in Euclidean
space and quantized through the functional integral as usual.
In this and the following section, dimensional regularization 
is employed.
The notational conventions are summarized
in appendix A.


\lussubsec 2.1 Flow equations

Let $A_{\mu}(x)$ be the $\SUn$ gauge potential,
$\psi(x)$ the quark field and $\psibar(x)$ the antiquark field
integrated over in the functional integral. The latter 
carry Dirac, colour and flavour indices, which are usually suppressed 
for simplicity. 

The Yang--Mills gradient flow evolves the gauge field as a function 
of a parameter $t\geq0$ that is referred to as the flow time.
Starting from the fundamental gauge field,
\equation{
  \left.B_{\mu}\right|_{t=0}=A_{\mu},
  \enum
}
the time-dependent field $B_{\mu}(t,x)$
is determined by the differential equation
\equation{
  \partial_tB_{\mu}=D_{\nu}G_{\nu\mu},
  \enum
  \nexteq{2ex}
  G_{\mu\nu}=\partial_{\mu}B_{\nu}-\partial_{\nu}B_{\mu}+[B_{\mu},B_{\nu}],
  \qquad
  D_{\mu}=\partial_{\mu}+[B_{\mu},\,\cdot\;]
  \enum
}
(see ref.~[\cite{WilsonFlow}] for an introduction to the subject).

As already mentioned in sect.~1, the extension of the flow to the
quark fields con\-sidered in this paper is a minimal one, where
the evolution of the gauge field is left unchanged.
The gauge field however appears in the flow 
equations\kern1pt\sfn{$\dagger$}{%
The left-action of any differential operator $\Delta$
(or of a difference operator in lattice field theory) 
is defined through the requirement that the relation
$\eta^{\dagger}\lvec{\Delta}=(\Delta\eta)^{\dagger}$
holds for all complex-valued fields $\eta$.}
\equation{
  \partial_t\chi=\Delta\chi,
  \qquad
  \partial_t\chibar=\chibar\lvec{\Delta},
  \enum
  \nexteq{2.5ex}
  \Delta=D_{\mu}D_{\mu},\qquad D_{\mu}=\partial_{\mu}+B_{\mu},
  \enum
}
which, together with the initial conditions
\equation{
  \left.\chi\right|_{t=0}=\psi,
  \qquad
  \left.\chibar\right|_{t=0}=\psibar,
  \enum
}
define the time-dependent quark and antiquark 
fields $\chi(t,x)$ and $\chibar(t,x)$.
In the flow equations (2.4), 
the gauge-covariant Laplacian $\Delta$ 
could be replaced by the square 
of the Dirac operator $\slash{D}$, for example, 
or be multiplied by a proportionality factor,
but such modifications
do not appear to offer any advantages.

The flow of quark fields introduced in this section 
is similar to the source 
smoothing operations proposed many years ago in
refs.~[\cite{Guesken},\cite{AlexandrouEtAl}].
With respect to these popular methods, there are, however, 
a few important differences,
one of them being the fact that the flow
operates in four dimensions rather than on the fields on an equal-time
hyper-plane. Moreover, the flow time varies continuously and the gauge field
is not set to the fundamental gauge field, but is
evolved together with the quark field. 


\lussubsec 2.2 Local composite fields

Since the flow equations are gauge-covariant, the time-dependent
fields transform in the same way under gauge transformations as
the fundamental fields. Examples of 
gauge-invariant composite fields
that are local both in space-time and in flow time are the densities
\equation{
  E_t(x)=-\frac{1}{2}\tr\{G_{\mu\nu}(t,x)G_{\mu\nu}(t,x)\},
  \enum
  \nexteq{2.5ex}
  S^{rs}_t(x)=\chibar_r(t,x)\chi_s(t,x),
  \qquad
  P^{rs}_t(x)=\chibar_r(t,x)\dirac{5}\chi_s(t,x),
  \enum
}
where $r,s$ are flavour indices.

Through the initial conditions (2.1),(2.6), 
the time-dependent fields depend on the 
fundamental fields. 
From the point of view of the QCD functional 
integral, composite fields like the densities (2.7),(2.8)
are therefore observables, similar to Wilson loops or the 
ordinary pseudo-scalar densities, for example.
In the following, the quantities of
interest are the correlation functions of 
these fields,
i.e.~the expectation values of products of local fields composed 
from the basic fields at any flow time. 


\lussubsec 2.3 Perturbation theory

As explained in refs.~[\cite{WilsonFlow},\cite{RenFlow}],
such correlation functions can be expanded in powers of the 
gauge coupling in a straightforward manner. 
First the flow equations are solved in powers of 
the fundamental fields. The substitution of 
these expansions in the local fields then leads to linear 
combinations
of correlation functions of the fundamental fields that
can be computed using the standard QCD Feynman rules.

In perturbation theory,
the flow equations (2.2) and (2.4) are replaced by
\equation{
  \partial_tB_{\mu}=D_{\nu}G_{\nu\mu}+\alpha_0D_{\mu}\partial_{\nu}B_{\nu},
  \enum
  \nexteq{2.5ex}
  \partial_t\chi=\Delta\chi-
  \alpha_0\partial_{\nu}B_{\nu}\chi,
  \qquad
  \partial_t\chibar=\chibar\lvec{\Delta}+
  \alpha_0\chibar\partial_{\nu}B_{\nu},
  \enum
}
in order to avoid some technical subtleties.
The terms proportional to the parameter $\alpha_0>0$ 
serve to damp the gauge modes
and could be removed by a time-dependent gauge transformation
[\cite{WilsonFlow}].
Correlation functions of gauge-invariant 
fields are therefore not affected by these extra terms.

The expansion of the time-dependent quark field $\chi(t,x)$ 
in powers of
the fundamental fields is obtained by iterating the 
quark flow equation in its integral form,
\equation{
  \chi(t,x)=\int\rmd^{D}\kern-1pt y\left\{K_t(x-y)\psi(y)
  +\int_0^t\rmd s\,K_{t-s}(x-y)\Delta'\chi(s,y)\right\},
  \enum
  \nexteq{3.0ex}
  \Delta'=(1-\alpha_0)\partial_{\nu}B_{\nu}
  +2B_{\nu}\partial_{\nu}+B_{\nu}B_{\nu},
  \enum
}
and the corresponding integral
equation for the time-dependent gauge field
[\cite{WilsonFlow}], where
\equation{
  K_t(z)={\rme^{-z^2/4t}\over(4\pi t)^{D/2}}
  \enum
}
denotes the heat kernel of the Laplacian in $D$ dimensions. 
At the lowest order in the gauge coupling, 
the interaction term in eq.~(2.11) can be dropped
and the expression
\equation{
  \langle\chi(t,x)\chibar(s,y)\rangle=
  \int{\rmd^{D}\kern-1pt p\over(2\pi)^D}\,
  \rme^{ip(x-y)}
  {\rme^{-(t+s)p^2}\over M_0+i\slash{p}}+\rmO(g_0^2)
  \enum
}
is then obtained for the two-point function of the time-dependent
quark field, where $g_0$ and $M_0$ are the bare coupling 
and (diagonal) quark mass matrix.
The smoothing character of the quark flow
is evident from these equations.
Moreover,
the smoothing range at leading order, $\sqrt{8t}$, 
is seen to be the same as in the case of the gauge field.

There are $8$ self-energy diagrams that contribute
to the two-point function (2.14) at one-loop order of perturbation theory.
The vertices in these diagrams derive from the QCD action
and from the iteration of the integral equation (2.11).
Following the lines of ref.~[\cite{RenFlow}], it is straightforward
to compute the ultraviolet divergent parts of the diagrams and
one then finds that the two-point function can be renormalized
by renormalizing the quark masses as usual and 
the fields according to
\equation{
  \chi=Z_{\chi}^{-1/2}\ren{\chi},
  \qquad
  \chibar=\ren{\chibar}Z_{\chi}^{-1/2}.
  \enum
}
In the $\MSbar$ scheme in $D=4-2\eps$ dimensions,
the calculation yields
\equation{
  Z_{\chi}=1+{3\CF\over16\pi^2\eps}g^2+\rmO(g^4),
  \qquad
  \CF={N^2-1\over2N},
  \enum
}
for the field renormalization constant,
$g$ being the renormalized coupling.


\lussubsec 2.4 Quark condensate at non-zero flow time

At non-zero flow times, the large momenta in the integral (2.14) 
are exponentially suppressed and the quark two-point function 
consequently has no singularities as
$(t,x)\to(s,y)$. The absence of short-distance singularities
is a general feature of the correlation functions at positive
flow times, which derives from the smoothing property of
the flow equations. 
In particular, the ``time-dependent quark condensates''
\equation{
  \Sigt^{rr}=-\langle S^{rr}_t(x)\rangle
  \enum
}
do not require additive renormalization.

In perturbation theory,
the renormalized condensates 
\equation{
  \rens{\Sigma}{t}^{rr}=Z_{\chi}\Sigt^{rr}
  \enum
}
can be worked out in powers of the renormalized coupling $g$,
with coefficients that depend on the renormalized quark masses 
$\mR{r},\mR{s},\ldots$ and the flow time $t$.
The leading-order term in four dimensions is given by
\equation{
  \left.\rens{\Sigma}{t}^{rr}\right|_{g=0}=
  {2N\mR{r}\over(4\pi)^2t}
  \int_0^{\infty}\rmd v\,{\rme^{-vz}\over(1+v)^2},
  \qquad z=2t\mR{r}^2,
  \enum
}
and thus vanishes in the chiral limit,
consistently with the fact that chiral symmetry can only
be spontaneously broken at the non-perturbative level.
